\chapter{Mirror Equation}

\section*{Introduction}

The mirror equation relates the object distance, image distance, and focal length for a spherical mirror. It is the foundation of geometric optics for curved reflecting surfaces, first systematically developed through contributions from Euclid (catoptrics, c.~300~BCE), Ibn al-Haytham (Optics, 1021), and Johannes Kepler (Dioptrice, 1611). The equation applies to both concave (converging) and convex (diverging) mirrors under the paraxial approximation, where rays are close to the optical axis.

The mirror equation encodes the law of reflection ($\theta_{i} = \theta_{r}$) for a spherical surface. For a concave mirror, parallel rays converge at the focal point; for a convex mirror, they appear to diverge from a virtual focal point behind the mirror. The equation $1/f = 1/d_{o} + 1/d_{i}$ is a direct consequence of similar triangles in the paraxial limit and governs how mirrors form images in telescopes, makeup mirrors, security mirrors, and reflecting telescopes.


\begin{equation}
\boxed{\; \frac{1}{f} = \frac{1}{d_{o}} + \frac{1}{d_{i}}\;}
\label{eq:mirror}
\end{equation}

The magnification is:
\begin{equation}
m = -\frac{d_{i}}{d_{o}}
\label{eq:magnification}
\end{equation}

\section*{Fundamental Equations Used}

The derivation starts from the geometry of a concave spherical mirror combined with the law of reflection.

\section*{Assumptions}

\textbf{Assumptions:} Paraxial rays (small angles); spherical mirror surface; single reflection; mirror is in vacuum or air.


\textbf{Limitations:} Spherical aberration for rays far from the axis; does not account for diffraction; limited to specular (not diffuse) reflection.


\section*{Mathematical Derivation}

\textbf{Step 1: Set up the geometry of a concave spherical mirror.}

Consider a concave spherical mirror of radius $R$ with center of curvature $C$. The focal length is $f = R/2$. Place the mirror vertex at the origin, the optical axis along the $x$-axis, and an object at distance $d_o$ to the left of the mirror. A ray from the top of the object (height $h_o$) strikes the mirror and reflects according to the law of reflection.

\textbf{Step 2: Trace a ray parallel to the axis.}

A ray from the top of the object traveling parallel to the axis hits the mirror at height $h$ and reflects through the focal point $F$ (at distance $f$ from the vertex). In the paraxial limit ($h \ll R$), the angle of incidence equals the angle of reflection, and this ray crosses the axis at $f$.

\textbf{Step 3: Trace a ray through the center of curvature.}

A ray from the top of the object passing through the center of curvature $C$ (at distance $R = 2f$ from the vertex) strikes the mirror normally and reflects back on itself. This ray and the parallel ray intersect at the image point.

\textbf{Step 4: Use similar triangles to relate the distances.}

From the parallel ray: the triangle with base $d_o$ and height $h_o$ is similar to the triangle with base $d_i$ and height $h_i$ (image height). The ray through $C$ gives:
\begin{equation}
\frac{h_o}{d_o} = \frac{h_i}{d_i} \qquad\text{(from similar triangles with the object and image triangles)}.
\end{equation}
From the focal-point crossing, the small triangle near the mirror (height $h_o$, base $d_o - f$) is similar to the triangle at the focal point (height $|h_i|$, base $d_i - f$):
\begin{equation}
\frac{h_o}{d_o - f} = \frac{|h_i|}{d_i - f}.
\end{equation}

\textbf{Step 5: Combine to derive the mirror equation.}

Using $h_i/h_o = -d_i/d_o$ (the minus sign indicates an inverted image for real images):
\begin{equation}
\frac{-d_i/d_o \cdot d_o}{d_i - f} = \frac{d_o}{d_o - f}.
\end{equation}
Simplifying:
\begin{equation}
\frac{-d_i}{d_i - f} = \frac{d_o}{d_o - f}.
\end{equation}
Cross-multiplying:
\begin{equation}
-d_i(d_o - f) = d_o(d_i - f) \quad\Longrightarrow\quad -d_i d_o + d_i f = d_o d_i - d_o f.
\end{equation}
Collecting terms:
\begin{equation}
2d_o d_i = f(d_o + d_i) \quad\Longrightarrow\quad \boxed{\frac{1}{f} = \frac{1}{d_o} + \frac{1}{d_i}}.
\end{equation}

\textbf{Step 6: Derive the magnification formula.}

From the similar-triangle relation:
\begin{equation}
m = \frac{h_i}{h_o} = -\frac{d_i}{d_o}.
\end{equation}
When $|m| > 1$, the image is enlarged; when $m < 0$, the image is inverted. For a plane mirror ($R \to \infty$, $f \to \infty$), the equation gives $d_i = -d_o$, confirming a virtual image at the same distance behind the mirror.

\section*{Characteristics of the Equation}

\begin{itemize}
  \item \textbf{Concave mirror:} $f > 0$; real images form when $d_{o} > f$; virtual images when $d_{o} < f$.
  \item \textbf{Convex mirror:} $f < 0$; always forms a virtual, upright, diminished image.
  \item \textbf{Object at infinity:} Image forms at the focal point ($d_{i} = f$).
  \item \textbf{Object at the focal point:} Image is at infinity (collimated beam).
  \item \textbf{Magnification:} $|m| > 1$ for enlarged images; $m < 0$ for inverted images; $m > 0$ for upright images.
\end{itemize}
\section*{Solution Methods}

\begin{enumerate}
  \item \textbf{From similar triangles:} Draw a ray from the top of the object through the center of curvature (reflects back on itself) and a ray parallel to the axis (reflects through the focal point). The intersection gives the image, and similar triangles yield Eq.~\eqref{eq:mirror}.
  \item \textbf{From Fermat's principle:} The optical path length from object to image via reflection on the mirror is stationary. For a parabolic mirror, all parallel rays have equal path lengths, giving perfect focusing; for a spherical mirror, the paraxial limit gives the same result approximately.
  \item \textbf{From the reflector equation:} A parabolic mirror ($z = r^{2}/4f$) exactly focuses parallel rays. The spherical mirror is a paraxial approximation to the paraboloid.
\end{enumerate}
\section*{Verifications}

The mirror equation is verified in every optics laboratory through image-distance measurements with benches and screens. The formation of real images at predicted positions and the measurement of magnification ratios agree with the equation to the limits set by spherical aberration and alignment precision.
\section*{Applications}

\begin{itemize}
  \item \textbf{Astronomical telescopes:} Newtonian and Cassegrain reflectors use concave primary mirrors to collect and focus starlight.
  \item \textbf{Solar concentrators:} Parabolic dishes focus sunlight to a point for thermal energy collection.
  \item \textbf{Dental and shaving mirrors:} Concave mirrors with $d_{o} < f$ produce magnified upright images.
  \item \textbf{Security and vehicle mirrors:} Convex mirrors provide wide-field views by forming diminished virtual images.
\end{itemize}
\section*{What Richard Feynman Would Have Asked}

\subsection*{Plain-English Translation}
\textit{``If I had to explain the mirror equation to someone who has never studied optics, how would I describe it?''}

If you hold a curved mirror and shine a flashlight at it, the reflected light forms an image somewhere. The mirror equation tells you exactly where: it connects three numbers---how far the object is, how far the image forms, and the curvature of the mirror. If the mirror is very curved (small $f$), the image is close; if the mirror is almost flat (large $f$), the image is far away. It is a bookkeeping rule for where light goes after reflection.

\subsection*{Localized Mechanics}
\textit{``At the level of a single ray reflecting from a small patch of the mirror, what determines the image location?''}

Consider a ray from the top of an object striking a point on the mirror at height $h$ above the axis. The law of reflection ($\theta_{i} = \theta_{r}$) determines the direction of the reflected ray. For a spherical mirror, the normal at the point of incidence points toward the center of curvature $C$. In the paraxial limit ($h \ll R$), the angle of incidence is approximately $h/R$ (measured from the axis), and the reflected ray's angle determines where it crosses the axis. Summing contributions from all such rays, they converge (or diverge) from a single image point---this is the content of the mirror equation. The paraxial approximation is the statement that $\sin\theta \approx \theta$, which is the same approximation that underlies the thin lens equation.

\subsection*{Testing Extreme Limits}
\textit{``What happens to the image as the object approaches the mirror, moves to infinity, or when the mirror radius becomes extreme?''}

As $d_{o} \to \infty$, $d_{i} \to f$: the image forms at the focal point regardless of where the object is, as long as it is very far away. This is why telescopes can image stars at any distance. As $d_{o} \to f^{+}$, $d_{i} \to \infty$: the image recedes to infinity, producing a collimated beam (this is how searchlights and flashlights work---place the bulb at the focus). As $d_{o} \to 0$, $d_{i} \to 0$: the image coincides with the object, with unit magnification. For a very flat mirror ($R \to \infty$, $f \to \infty$), the equation gives $d_{i} = -d_{o}$, recovering the plane mirror result. For a very curved mirror ($R \to 0$), $f \to 0$, and the image always forms very close to the mirror surface.

\subsection*{Dimensionality and Geometry}
\textit{``What is the geometric structure behind the mirror equation, and how does it relate to the shape of the mirror?''}

The mirror equation is a projective statement about conic sections. A parabolic mirror ($z = r^{2}/4f$) exactly maps parallel rays to a single point---this is a geometric property of parabolas known since Apollonius. A spherical mirror is an approximation to a paraboloid, valid when $r \ll R$. The image location is the intersection of the axis with the reflected wavefront, which, for a spherical mirror, is a sphere centered on $C$ and passing through the image point. The equation $1/f = 1/d_{o} + 1/d_{i}$ is equivalent to saying that the object and image distances are conjugate under inversion through the focal length---a transformation related to the Cayley--Klein model of hyperbolic geometry. The deeper structure is that of conformal mapping: the mirror transforms spherical wavefronts from the object into spherical wavefronts from the image, preserving angles.

\subsection*{Cross-Disciplinary Connection}
\textit{``Where does the structure of the mirror equation---relating distances through a reciprocal sum---appear outside of optics?''}

The form $1/f = 1/d_{o} + 1/d_{i}$ appears throughout physics wherever a quadratic potential maps inputs to outputs. In electrostatics, the lens maker's equation for thin lenses has the identical form. In gravitation, the two-body problem can be recast in terms of an ``effective focal length'' of the Keplerian orbit. In electronics, the formula $1/R_{\text{eq}} = 1/R_{1} + 1/R_{2}$ for parallel resistors has the same mathematical structure. In economics, the ``gravity model'' of trade between cities uses reciprocal-distance relationships. The unifying structure is that of the harmonic mean, which appears whenever quantities combine in parallel or when inversion is a natural symmetry of the problem. In optics specifically, the mirror equation generalizes to the thick-mirror equation, the matrix method (ABCD matrices), and ultimately to the wave-optical diffraction integral.

% ============================================================
% CHAPTER 43 — Navier–Stokes Equation
% ============================================================
\section*{FAQ}

\begin{enumerate}
\item \textbf{Why doesn't a spherical mirror focus perfectly?}
A sphere is not a paraboloid. Rays far from the axis are focused closer to the mirror than paraxial rays, causing spherical aberration. A parabolic mirror has no spherical aberration for on-axis sources but introduces coma for off-axis sources.

\item \textbf{Can a convex mirror form a real image?}
No, not for a single mirror with a real object. The image is always virtual, upright, and diminished. However, in combination with other optical elements, a convex mirror can contribute to a system that forms real images.

\item \textbf{How does the mirror equation change for a plane mirror?}
A plane mirror has $R = \infty$ and $f = \infty$. The equation gives $d_{i} = -d_{o}$, meaning the image is virtual, at the same distance behind the mirror as the object is in front.
\end{enumerate}
\section*{Important Bibliography}

\begin{enumerate}
\item Hecht, E., \textit{Optics}, 5th ed. (Pearson, 2017), Chapter 5.
\item Jenkins, F.A. and White, H.E., \textit{Fundamentals of Optics}, 4th ed. (McGraw-Hill, 1976), Chapter 6.
\item Serway, R.A. and Jewett, J.W., \textit{Physics for Scientists and Engineers}, 10th ed. (Cengage, 2019), Chapter 35.
\item Born, M. and Wolf, E., \textit{Principles of Optics}, 7th ed. (Cambridge University Press, 1999), Chapter 4.
\end{enumerate}

